% page 541 du fac-simile (image p-540 du scan)
% bloc 185.4 x 233.3 mm ; marge gauche 16.6 mm
\pgc{-5.080mm}{0.000mm}{
\l{\cel{62}533}
\l{}
\l{\cel{5}\sou{sismometro}.\cel{2}Sinonimo di\cel{1}\sou{sismografo}.DEFIRS.}
\l{}
\l{\cel{5}\sou{sismoskopo}. Sinonimo di sismografo.}
\l{}
\l{\cel{5}\sou{sistemo}. I. Ensemblo di qua la parti esas koordinita segun}
\l{\cel{9}lego) (aludante enti, institucuri, doktrini, e c.) - II.}
\l{\cel{8}(\sou{anat.}) Ensemblo ek la organi qui esas uzata por ta o ca}
\l{\cel{8}funciono di la ekonomio che la animalo. (Kom ex.: la sistemo}
\l{\cel{8}digestala, respirala). - DEFIRS.}
\l{}
\l{\cel{5}\sou{sistematiko}. Parto di la natur-historio qua koncernas la klasifiko}
\l{\cel{8}di la enti, e qua inkluzas la deskripto di olci. - DEFIS.}
\l{}
\l{\cel{5}\sou{sistolo}. (\sou{anat.})\cel{2}Movo per qua la fibri muskulala di la kordio}
\l{\cel{8}kontraktesas por pulsar sango aden la arterii. - DEFIS.}
\l{}
\l{\cel{5}\sou{sistro.}\cel{2}(\sou{antique}) Muzik-instrumento di la Egiptiani antiqua,}
\l{\cel{8}ek mikra ringego de metalo trairata da vergeti qui sonoris}
\l{\cel{8}kande onu agitis li.}
\l{}
\l{\cel{5}\sou{sitelo}. Vazo cilindra, ek ligno, zinko, e c., kun anso, kapaca de}
\l{\cel{8}plura litri, qua uzesas por cherpar, transportar,aquo o liquido}
\l{\cel{8}altra. -\cel{2}I.}
\l{}
\l{\cel{5}\sou{situar}. (\sou{trans}.)Pozar en ta o ca loko (pri expozeso, aspekto, pozi-}
\l{\cel{8}ciono defensala, e c.).DEFIS.}
\l{}
\l{\cel{5}\sou{sivo}. (\sou{tekn.}) Cirklo-kurvo de ligno sur qua esas tensita telo kun}
\l{\cel{8}mashi plu o min densa,uzata por kriblagar materii pulveratra}
\l{\cel{8}(farino, gipso, e c.) - DE.}
\l{}
\l{\cel{5}\sou{sive}. (\sou{gram.}) Konjunciono qua indikas alternativo, uzata avan verbo}
\l{\cel{8}expresita o tacita. L.}
\l{}
\l{\cel{5}\sou{sizar}. (\sou{trans.}) Prenar, manuo-kaptar, rapide e vigoroze (ulo, ulu).}
\l{\cel{7}- EF.}
\l{}
\l{\cel{5}\sou{sizigio}. (\sou{astron.})\cel{2}Situeso di Suno e diLuno, kande ta du astri esas}
\l{\cel{8}od en kunjunto fnova Luno\textquotedbl{}), od en opozeso (\textquotedbl{}plena Lunof).-}
\l{\cel{8}DEFIRS.}
\l{}
\l{\cel{5}\sou{skabelo}. I. Sidilo ek ligno, poke alta, sen brakii nek dors-apogilo.}
\l{\cel{8}- II. Quaza eskalereto movebla, kun du o tri gradi, sur qua onu}
\l{\cel{8}acensas por atingar ulo. - FIS.}
\l{}
\l{\cel{5}\sou{skabio}. (\sou{patol.}) Morbo kontagiala di la pelo, quan produktas la}
\l{\cel{8}prezenteso di la insekto \textquotedbl{}skaro\textquotedbl{} (L.\cel{1}\sou{acarus}). EF.}
\l{}
\l{\cel{5}\sou{skabino}. Judiciisto municipala di ula urbi. - FI.}
\l{}
\l{\cel{5}\sou{skabio}so. (\sou{bot.})\cel{2}Planto herbatra, di qua la flori esas violea, pur-}
\l{\cel{8}purea od (en ula kazi) blanka. - L.\cel{1}\sou{scabiosa.}\cel{2}DEFIRS.}
}
